\section{Localized losses and the Gaussian posterior}
\label{sec:localized-proof}
This section uses the common design-norm event in Section~\ref{sec:sharp-design}.
The statements below use bounded labels $|Y|\le1$ and a density on
$[-1,1]$ bounded above and below by positive constants. Put
$f_0(x)=\mathbb E[Y\mid X=x]$, $A(f)=P_X(f-f_0)^2$,
and $\xi=Y-f_0(X)$. In particular $|\xi|\le2$.

\paragraph{A self-contained entropy bound.}
For $P\ge1$, let $\mathcal V_P$ contain continuous functions with
supremum, Lipschitz constant and total variation of their distributional
second derivative all at most $cP$, where $c$ is an absolute constant.
Then
\[
 \log N_\infty(\varepsilon,\mathcal V_P)
 \le C\sqrt{P/\varepsilon}\log(CP/\varepsilon),
 \qquad 0<\varepsilon\le P. \tag{E}
\]
Here and below an extra universal logarithm may be inserted without
changing the conclusions. To see (E), partition the interval both by
a uniform grid with $k$ bins and by quantiles of the second-derivative
variation measure with at most $k$ internal cuts. The common
refinement has at most $2k+1$ intervals. On each open interval its
length is at most $2/k$ and its variation measure is at most $cP/k$;
atoms at endpoints do not affect linear interpolation. The linear
interpolant therefore differs uniformly from the original function
by at most $CP/k^2$. Choose $k=C\sqrt{P/\varepsilon}$.
Snap the interpolant's knots to an $\varepsilon/(CP)$ grid and
its values to an $\varepsilon/C$ grid. Its Lipschitz constant is
at most $cP$, so snapping and merging colliding knots changes the
uniform error by at most $C\varepsilon$. The number of resulting
descriptions proves (E). The affine part is included by the
supremum and Lipschitz bounds.

\paragraph{Two multiplier estimates.}
Fix a deterministic bounded comparator $g$, with $|g|\le1$ and
path variation $P_g$. For any shallow ReLU output $f$, write
$h=f-g$ and $P=1+P_f+P_g$. Both $h$ and the possible values of $f$
belong to a constant enlargement of $\mathcal V_P$.
For $0<\eta<\eta_0$, where $\eta_0>0$ depends only on the two density bounds, with probability at least $1-\delta$,
simultaneously over all $f$,
\begin{align}
 |P_n\xi h|
 &\le \eta P_nh^2+
 C\eta^{-3/5}P^{2/5}n^{-4/5}\mathcal L^C
 +C\eta^{-1}n^{-1}\mathcal L^C, \tag{M1}\\
 |(P_n-P_X)((f_0-g)h)|
 &\le \eta P_Xh^2+
 C\eta^{-3/5}P^{2/5}n^{-4/5}\mathcal L^C
 +CPn^{-13/8}\mathcal L^C\\
 &\quad+C\eta^{-1/2}P^{1/2}n^{-3/2}\mathcal L^C
 +C\eta^{-1}n^{-1}\mathcal L^C. \tag{M2}
\end{align}
Here $\mathcal L=1+\log(n/\delta)+\log(P+1)$; one can replace its
$P$ dependence on dyadic shells by the corresponding shell logarithm.
Both events depend only on the function classes, not on learner
width or Gaussian scale.

For (M1), condition on the inputs. The process $P_n\xi h$ is
sub-Gaussian for the empirical $L^2$ metric, with variance proxy
$C/n$. On $P_nh^2\le r^2$, Dudley's chaining bound and (E) give
\[
 \mathbb E_\xi\sup |P_n\xi h|
 \le C n^{-1/2}P^{1/4}r^{3/4}\mathcal L^C.
\]
The usual sub-Gaussian chaining tail adds
$Cr\sqrt{t/n}$. Peel in $P$ and $r$ and use Young's inequality.
Terminate the radius peeling at $r=1/n$. For the smaller terminal
class, $|P_n\xi h|\le2r\le2/n$ directly. For the design
multiplier use the localized estimate at $r=1/n$. Thus no infinite
peeling or uncontrolled small-radius limit is needed.

For (M2), localize first to $P_Xh^2\le r^2$. Apply the design-norm
lemma at a fixed tolerance, say $1/4$, and failure $n^{-10}\delta$.
It bounds all empirical radii by
\[
 C r+C P\{\log(Cn/\delta)/n\}^{3/2}.
\]
Conditional Rademacher chaining using (E), with the bounded
multiplier $f_0-g$, gives expected supremum at most
\[
 C n^{-1/2}P^{1/4}r^{3/4}\mathcal L^C
       +CPn^{-13/8}\mathcal L^C.
\]
The exceptional design event contributes at most $CPn^{-10}\delta$.
The exponent $13/8$ follows by inserting the additional empirical
radius $P n^{-3/2}$ in $n^{-1/2}P^{1/4}r^{3/4}$.

The lower density bound supplies a localized envelope:
\[
 \|h\|_\infty\le C\{r+P^{1/3}r^{2/3}\}.
\]
Indeed, if $|h(x_0)|=A$, a one-sided interval of length
$c\min(1,A/P)$ has $|h|\ge A/2$. Hence
$r^2\ge c A^2\min(1,A/P)$, proving the two alternatives.
Bousquet's inequality therefore adds
\[
 Cr\sqrt{t/n}+C(r+P^{1/3}r^{2/3})t/n,
\]
using variance at most $4r^2$. Its square-root term involving the
expected supremum is absorbed by that expectation and the displayed
envelope term. The population localization makes this envelope
deterministic on the whole class. Young's inequality bounds
$P^{1/3}r^{2/3}/n$ by
$\eta r^2+C\eta^{-1/2}P^{1/2}n^{-3/2}$.
Dyadic peeling and the same terminal radius $1/n$ prove (M2).
No independence of $f_0(X)-g(X)$ and $X$ is assumed.

\paragraph{A relative squared-loss inequality.}
Let
\[
 \widehat R(f)=P_n(f-Y)^2,\qquad
 D_n=C\eta^{-5}\{\log(Cn/\delta)/n\}^{3}.
\]
Combine (M1), (M2), and the two directions of the design-norm lemma:
\begin{equation}
 \begin{aligned}
 |(\widehat R(f)-\widehat R(g))-(A(f)-A(g))|
 \le{}& C\eta\{A(f)+A(g)\}
 +C\eta^{-3/5}P^{2/5}n^{-4/5}\mathcal L^C\\
 &+CPn^{-13/8}\mathcal L^C
 +C\eta^{-1/2}P^{1/2}n^{-3/2}\mathcal L^C\\
 &+C\eta^{-1}n^{-1}\mathcal L^C+D_nP^2.
 \end{aligned}\tag{L}
\end{equation}
Indeed the empirical/population quadratic-norm difference for
$h$ is at most $C\eta P_Xh^2+D_nP^2$, and
$P_Xh^2\le2A(f)+2A(g)$. The multiplier (M1) initially involves
$P_nh^2$, which the reverse norm inequality controls by the same
quantity. Constants in the fine-tolerance design lemma can be
absorbed by replacing its input tolerance by a fixed multiple of
$\eta$.

\paragraph{The actual Gaussian posterior.}
Use the original shallow independent Gaussian prior
\[
 a_j,b_j\sim N(0,\nu),\quad v_j\sim N(0,\nu/N^2),
 \quad \omega=N/\nu,\quad 0<\nu\le1,
\]
and define
\[
 Q=N^{-1}\sum(a_j^2+b_j^2)+N\sum v_j^2,\quad
 \Pi_\lambda(dw)\propto
 e^{-\omega\widehat R(z_w)/(2\lambda)}\mu(dw).
\]
Then $P(z_w)\le Q/2$, $\omega Q\sim\chi^2_{3N}$, and the physical
Gaussian observation variance is $n\lambda\nu/N$.
The Gibbs identity and the Gaussian exponential moment imply
\[
 \Pi_\lambda[\widehat R(z)+(\lambda/2)Q]
 \le F_\lambda+3\lambda\log2,\qquad
 F_\lambda=-\frac{2\lambda}{\omega}\log Z_\lambda. \tag{G}
\]
In fact
$\operatorname{KL}(\Pi_\lambda\Vert\mu)
 \ge\omega\Pi_\lambda Q/4-(3N/2)\log2$ and $N/\omega\le1$.

For $0<\lambda\le1$ there is also the deterministic bound
\[
 \Pi_\lambda Q^2\le C(1+\lambda^{-2}). \tag{T}
\]
The prior empirical risk is at most two for every unit-bounded
sample, so Jensen gives $F_\lambda\le2$. A chi-square Chernoff
bound gives $\mu(Q>u)\le e^{-c\omega u}$ above a fixed universal
threshold. Since $Z_\lambda\ge e^{-\omega/\lambda}$,
the posterior tail is at most
$\min(1,e^{\omega/\lambda-c\omega u})$; integrate it and use
$\omega\ge1$ to obtain (T).

For the following uniform statistical bound assume $P_g\le n^4$.
Truncate posterior energy at $Q\le n^{20}/\lambda$; the tail estimate
used in (T) bounds every omitted risk or quadratic-energy integral
by $O(n^{-10})$. On the retained range all shell logarithms are
$O(\log(n/\delta))$ when $n^{-2}\le\lambda\le1$.
For any $P\ge1$ and $a>0$,
\[
 C\eta^{-3/5}P^{2/5}n^{-4/5}\mathcal L^C
 \le a\lambda P+
 C_a\eta^{-1}n^{-4/3}\lambda^{-2/3}\mathcal L^C. \tag{Y}
\]
For the additional sublinear term in (M2), Young gives
\[
 C\eta^{-1/2}P^{1/2}n^{-3/2}\mathcal L^C
 \le a\lambda P+C_a\eta^{-1}\lambda^{-1}n^{-3}\mathcal L^C.
\]
The linear term $Pn^{-13/8}$ is absorbed provided
\[
 \lambda\ge C n^{-13/8}\log^k(n/\delta)
\]
for one sufficiently large fixed $k$. The Gaussian tail truncation
above makes every retained shell logarithm $O(\log(n/\delta))$;
the discarded posterior risk is negligible uniformly in $N,\nu$.

A balanced Gaussian comparison around $g$ gives, for every
$N$ above the width of $g$,
\[
 F_\lambda\le\widehat R(g)
       +C/n+\lambda C[P_g+\log(n(P_g+1))].
\]
Choose $a$ in (Y) small enough that the resulting coefficient
of $Q$ in (L) is at most $\lambda/2$, and integrate (L).
Using (G), (T), and the comparison gives
\begin{align}
 \Pi_\lambda A(z)
 \le &(1+C\eta)A(g)
 +C\lambda[P_g+\log(n(P_g+1))]
 +C\eta^{-1}n^{-4/3}\lambda^{-2/3}\mathcal L_n^C\\
 &+C\eta^{-1}\lambda^{-1}n^{-3}\mathcal L_n^C
 +C\eta^{-5}n^{-3}\mathcal L_n^C
       [\,\lambda^{-2}+(P_g+1)^2\,]+C\eta^{-1}n^{-1}\mathcal L_n^C, \tag{O}
\end{align}
where $\mathcal L_n$ contains only logarithms in
$n,1/\delta,P_g+1$. Jensen gives the same bound for the ordinary
raw posterior mean. All events are shared by every eligible
finite width and every $0<\nu\le1$.

\paragraph{A universal comparison-width guarantee.}
Take $\lambda=n^{-3/2}\log^K(n/\delta)$ with one sufficiently large
fixed $K$, and put $\eta=1/\log n$. Restrict to the nontrivial
comparison widths $(C+1)\sqrt\lambda\le1$.
The bounded spline compiler applied to a clipped width-$C$ univariate
shallow ReLU function has width at most $6C+6$, risk increment
$O(Ct)$ and path variation $P_g\le C'(C+1)/t$.
Use $t=\sqrt\lambda$ in (O). The compiler and comparison penalty
are $\widetilde O((C+1)n^{-3/4})$. The noise remainder is
$\widetilde O(n^{-1/3})$, and the new envelope remainder is
$\widetilde O(n^{-3/2})$.
The posterior design term is
\[
 C\eta^{-5}n^{-3}\mathcal L_n^c\lambda^{-2}
 \le C\log^{c+5-2K}(n/\delta)\le C/\log n
\]
after increasing $K$. Its comparator-energy term is this quantity
times $O((C+1)^2\lambda)$, and is therefore no larger in the
stated range. Thus
\[
 R(\bar z_\lambda)\le\Pi_\lambda R(z)
 \le R_0+\{1+O(1/\log n)\}
       \inf_{g\in\mathcal G_C}\{R(g)-R_0\}
 +\widetilde O((C+1)n^{-3/4}+n^{-1/3})+O(1/\log n).
 \tag{C}
\]
Clipping a proof comparator to $[-1,1]$ cannot increase its squared
risk for unit-bounded labels, and creates at most $3C+2$ breakpoints.
The actual posterior remains raw and unchanged.
The same deterministic schedule works for all qualifying $C$.
A summable confidence allocation over integer comparison widths is
a proof union, not a prior over widths. In particular it permits
every fixed growth $C_n=n^u$ with $u<3/4$.
The logarithmic additive term is explicit: this universal schedule
is not the polynomial-accuracy bound of the declared-budget schedule
below. The conditions on $\eta,\lambda,P_g$ hold for sufficiently
large $n$; a bounded-risk fallback handles smaller samples.

\paragraph{Two coefficient-one schedules.}
The sharper design lemma yields two genuine coefficient-one
consequences, not only the coefficient-$1+o(1)$ interpretation.
All the following displays suppress fixed powers of logarithms;
these can be realized by enlarging the indicated tolerances and
likelihood scales by fixed logarithmic factors.

First take the universal schedule
\[
 \lambda=n^{-3/5},\qquad \eta=c n^{-3/10},
 \qquad t=\sqrt\lambda=n^{-3/10}.
\]
The likelihood schedule does not use $C$. Since $A(g)\le4$ for the
bounded proof comparator, its multiplicative slack is absorbed
additively. Substituting in (O) gives, for every $1\le C\le n$,
\[
 R(\bar z_\lambda)\le\inf_{g\in\mathcal G_C}R(g)
 +\widetilde O\left((C+1)n^{-3/10}
                   +(C+1)^2n^{-9/10}\right). \tag{U}
\]
The multiplier remainder has order $n^{-19/30}$ and the direct
$n^{-1}\eta^{-1}$ remainder has order $n^{-7/10}$; both are smaller
than the first displayed term. The design remainder from the
posterior energy is $n^{-3/10}$, and its comparator-energy term is
$(C+1)^2n^{-9/10}$. The hypotheses $\lambda\ge\widetilde C n^{-13/8}$ and
$n\eta^2\ge C\log(Cn/\delta)$ hold for sufficiently large $n$.
The constants depend on the two density bounds.

Second, if a comparison-width budget $C$ is declared, put
\[
 \varepsilon_{n,C}
 =\widetilde C(C+1)^{2/5}n^{-3/10},\qquad
 \eta=c\varepsilon_{n,C},\quad
 \lambda=c'\left(\frac{\varepsilon_{n,C}}{C+1}\right)^2,\quad
 t=c''\frac{\varepsilon_{n,C}}{C+1}.
\]
Choose the fixed constants and logarithmic powers in this order
so that the inequalities used above hold. For the nontrivial regime
$\varepsilon_{n,C}\le\varepsilon_0$, this yields
\[
 R(\bar z_\lambda)\le\inf_{g\in\mathcal G_C}R(g)
 +\widetilde O\left((C+1)^{2/5}n^{-3/10}\right). \tag{B}
\]
The learner uses $n,C$ and the confidence allocation to choose
one likelihood variance; it does not mix widths, priors or
temperatures. It still uses the same prior and any prescribed
finite learner width $N\ge6C+6$.

For the algebra, the compiler and penalty terms are $O(\varepsilon)$.
The strongest design requirement is
\[
 (C+1)^4 n^{-3}\varepsilon^{-9}\le C\varepsilon,
\]
which is exactly $\varepsilon\gtrsim(C+1)^{2/5}n^{-3/10}$.
The new envelope remainder requires only
$\varepsilon\gtrsim(C+1)^{1/2}n^{-3/4}$ and is weaker.
The noise multiplier asks only for
$\varepsilon\gtrsim(C+1)^{2/5}n^{-2/5}$, also weaker.
The comparator-energy design term is smaller than the posterior
one by a factor $O(\varepsilon^2)$.
In the nontrivial regime $\varepsilon_{n,C}\le\varepsilon_0$,
one has $C+1\lesssim n^{3/4}$ and hence
$\lambda\gtrsim n^{-3/2}$ with logarithmic slack. This exceeds
$n^{-13/8}\log^k(n/\delta)$ for sufficiently large $n$.
The other conditions $P_g\le n^4$ and
$n\eta^2\ge C\log(Cn/\delta)$ also hold.
Thus a sufficient sample size for excess $\varepsilon$ is
\[
 n\ge\widetilde O\left(
 (C+1)^{4/3}\varepsilon^{-10/3}\right). \tag{SC}
\]
Consequently every prescribed comparison width $C_n=n^u$, $u<3/4$,
has a vanishing power-rate remainder, at the deterministic scale
$\lambda_n=\widetilde\Theta(n^{-3/5-6u/5})$.
